\chapter{1964:Dorothy Crowfoot Hodgkin}

\label{chap:chemistry1964}

The Nobel Prize in Chemistry for the year 1964 was awarded to Dorothy Crowfoot Hodgkin.

\textbf{Motivation:}

for her determinations by X-ray techniques of the structures of important biochemical substances

\section{Dorothy Crowfoot Hodgkin}

\subsection*{Early Life and Education}

Dorothy Crowfoot Hodgkin was born on 1910-05-12 in Cairo, Egypt.

\subsection*{Career and Major Research}

Hodgkin studied chemistry at Somerville College, Oxford, and took her doctorate at the University of Cambridge in 1937 under J.D. Bernal. She returned to Oxford, where she spent her career, becoming Wolfson Research Professor of the Royal Society. She applied X-ray crystallography to increasingly complex biochemical substances, determining the structure of cholesterol in 1937, of penicillin in 1945, and of vitamin B12 in 1955.

\subsection*{Nobel Prize-winning Work}

The work for which Dorothy Crowfoot Hodgkin was awarded the Nobel Prize in Chemistry in 1964 was described by the Nobel Committee as: "for her determinations by X-ray techniques of the structures of important biochemical substances".

Hodgkin's determination of the structure of penicillin confirmed the previously disputed beta-lactam ring, and her solution of the structure of vitamin B12 revealed the large corrin ring and its central cobalt atom. Both structures were among the most complex molecules ever solved by X-ray analysis at the time.

\subsection*{Significance and Impact}

Hodgkin's structures of penicillin and vitamin B12 demonstrated that X-ray crystallography could solve molecules of biochemical and pharmaceutical importance, guiding later work on the synthesis and medical use of these substances. Her achievement inspired the use of the method throughout biochemistry.

\subsection*{Later Career and Honors}

Hodgkin continued X-ray analysis of insulin, work begun in the 1930s, and her group finally established the full insulin structure in 1969. She received many honors, including the Order of Merit in 1965. She died on 1994-07-29 in Shipston-on-Stour, England.

\subsection*{Legacy}

Hodgkin was one of the pioneers of macromolecular crystallography and one of the few women to receive the Nobel Prize in Chemistry. Her structures of cholesterol, penicillin, vitamin B12, and insulin remain landmarks of structural biology.

\subsection*{Selected Publications}

\begin{itemize}

\item \emph{Structure of Vitamin B12}, Nature, 1955.

\end{itemize}

\clearpage

\section*{Chapter Summary}

The 1964 prize recognized determinations of the structures of important biochemical substances by X-ray techniques. Dorothy Crowfoot Hodgkin solved the structures of cholesterol, penicillin, and vitamin B12 at Oxford, and her later work established the structure of insulin.

